\begin{table}[!htbp]\centering
\fontsize{9}{11}\selectfont\setlength{\tabcolsep}{4pt}
\caption{Barrido exploratorio ampliado: medias equiponderadas de seis datasets, con 47 condiciones por tamaño y clasificador.}\label{tab:tm-curva}
\begin{tabular}{@{}lrrr@{}}\toprule
k & Macro-F1 SVM & Macro-F1 ResMLP & \% dimensiones SVM\\\midrule
1 & 0,917448 & 0,917119 & 4,28\\
2 & 0,929892 & 0,932352 & 8,37\\
3 & 0,933921 & 0,933727 & 11,70\\
4 & 0,936554 & 0,933077 & 15,42\\
5 & 0,934857 & 0,935195 & 19,02\\
10 & 0,936764 & 0,933165 & 45,17\\
21 & 0,935616 & 0,929877 & 100,00\\
\bottomrule\end{tabular}
\par\smallskip\begin{minipage}{0.98\textwidth}\footnotesize La referencia dimensional son los 21 bloques del barrido, sin N-gramas. El máximo observado no es un tamaño óptimo confirmado.\end{minipage}
\end{table}
